% ===============================================================
% ML-05b: Musterlösung - Reguläre Verben auf -er/-ir
% Abgeleitet aus LE-05b (Abschnitt 9: Lösungen)
% ===============================================================

\documentclass[11pt, a4paper]{article}

\usepackage[utf8]{inputenc}
\usepackage[T1]{fontenc}
\usepackage[ngerman]{babel}
\usepackage{helvet}
\renewcommand{\familydefault}{\sfdefault}

\usepackage[margin=2cm]{geometry}
\usepackage{enumitem}
\usepackage{tabularx}
\usepackage{booktabs}

\usepackage{xcolor}
\usepackage{tcolorbox}

\usepackage{fancyhdr}
\usepackage{lastpage}

% Farben
\definecolor{spanisch}{HTML}{c41e3a}
\definecolor{grau}{HTML}{888888}
\definecolor{hellgrau}{HTML}{f5f5f5}
\definecolor{orange}{HTML}{b35c00}
\definecolor{loesung}{HTML}{c62828}
\definecolor{blau}{HTML}{2c5aa0}
\definecolor{gruen}{HTML}{2a7a4b}

\newcommand{\fachfarbe}{spanisch}

% Variablen
\newcommand{\thema}{Reguläre Verben auf -er/-ir}
\newcommand{\sequenz}{05}
\newcommand{\block}{b}
\newcommand{\fach}{Spanisch}
\newcommand{\klasse}{7}
\newcommand{\maxpunkte}{20}

% Kopf-/Fußzeile
\pagestyle{fancy}
\fancyhf{}
\renewcommand{\headrulewidth}{0.4pt}
\renewcommand{\footrulewidth}{0.4pt}
\lhead{\fach{} -- Klasse \klasse}
\chead{\textcolor{loesung}{\textbf{ML-\sequenz\block{} (MUSTERLÖSUNG)}}}
\rhead{}
\lfoot{}
\rfoot{Seite \thepage\ von \pageref{LastPage}}

% Befehle
\newcommand{\punktebox}[1]{\hfill\fbox{\small \_/#1}}
\newcommand{\lsg}[1]{\textcolor{loesung}{\textbf{#1}}}
\newcommand{\luecke}[1]{\rule{#1}{0.4pt}}
\newcommand{\es}[1]{\textcolor{\fachfarbe}{\textbf{#1}}}
\newcommand{\er}[1]{\textcolor{blau}{\textbf{#1}}}
\newcommand{\ir}[1]{\textcolor{gruen}{\textbf{#1}}}

% Merkkasten
\newtcolorbox{merkkasten}[1][]{
    colback=hellgrau,
    colframe=\fachfarbe,
    boxrule=1pt,
    arc=0pt,
    left=8pt, right=8pt, top=8pt, bottom=8pt,
    fonttitle=\bfseries\color{\fachfarbe},
    title=#1
}

% Hinweis
\newtcolorbox{hinweis}{
    colback=white,
    colframe=orange,
    boxrule=1pt,
    arc=0pt,
    left=5pt, right=5pt, top=5pt, bottom=5pt,
    fonttitle=\bfseries,
    title=Tipp
}

% Aufgabe
\newenvironment{aufgabe}[1]{%
    \vspace{10pt}
    \noindent\textbf{\textcolor{\fachfarbe}{Aufgabe #1}}
    \vspace{5pt}
    \begin{list}{}{\leftmargin=0pt}
    \item[]
}{%
    \end{list}
}

% ===============================================================
\begin{document}

% Kopfbereich
\begin{center}
    {\Large\bfseries\textcolor{\fachfarbe}{Arbeitsblatt: \thema}}\\[0.3em]
    {\large\textcolor{loesung}{\textbf{--- MUSTERLÖSUNG ---}}}
\end{center}

\vspace{5pt}

\noindent
\begin{tabularx}{\textwidth}{|X|X|X|}
\hline
\textbf{Name:} \textcolor{loesung}{Musterschüler/in} &
\textbf{Klasse:} \klasse &
\textbf{Datum:} \textcolor{loesung}{XX.XX.XXXX} \\
\hline
\end{tabularx}

\vspace{15pt}

% ===============================================================
% MERKKASTEN 1 (mit Lösungen)
% ===============================================================

\begin{merkkasten}[Merkkasten 1: Verben auf \er{-er} (z.B. comer{,} beber{,} leer{,} correr)]

\begin{center}
\begin{tabular}{|l|l|l|}
\hline
\textbf{Person} & \textbf{Pronomen} & \textbf{comer} (essen) \\
\hline
1. Sg. & yo & \lsg{como} \\
\hline
2. Sg. & tú & \lsg{comes} \\
\hline
3. Sg. & él/ella & \lsg{come} \\
\hline
1. Pl. & nosotros/-as & \lsg{comemos} \\
\hline
2. Pl. & vosotros/-as & \lsg{coméis} \\
\hline
3. Pl. & ellos/ellas & \lsg{comen} \\
\hline
\end{tabular}
\end{center}

\vspace{0.3em}
\textbf{Endungen -er:} \er{-o, -es, -e, -emos, -éis, -en}
\end{merkkasten}

\vspace{10pt}

% ===============================================================
% MERKKASTEN 2 (mit Lösungen)
% ===============================================================

\begin{merkkasten}[Merkkasten 2: Verben auf \ir{-ir} (z.B. vivir{,} escribir{,} abrir)]

\begin{center}
\begin{tabular}{|l|l|l|}
\hline
\textbf{Person} & \textbf{Pronomen} & \textbf{vivir} (leben/wohnen) \\
\hline
1. Sg. & yo & \lsg{vivo} \\
\hline
2. Sg. & tú & \lsg{vives} \\
\hline
3. Sg. & él/ella & \lsg{vive} \\
\hline
1. Pl. & nosotros/-as & \lsg{vivimos} \\
\hline
2. Pl. & vosotros/-as & \lsg{vivís} \\
\hline
3. Pl. & ellos/ellas & \lsg{viven} \\
\hline
\end{tabular}
\end{center}

\vspace{0.3em}
\textbf{Endungen -ir:} \ir{-o, -es, -e, -imos, -ís, -en}
\end{merkkasten}

\vspace{15pt}

% ===============================================================
% AUFGABE 1 (mit Lösungen)
% ===============================================================

\begin{aufgabe}{1}
Setze die richtige Form des Verbs in Klammern ein. \punktebox{5}
\end{aufgabe}

\begin{enumerate}[label=\alph*)]
    \item Yo \lsg{como} pizza. (comer)
    \item Tú \lsg{bebes} agua. (beber)
    \item María \lsg{lee} un libro. (leer)
    \item Nosotros \lsg{corremos} en el parque. (correr)
    \item Ellos \lsg{beben} mucha leche. (beber)
\end{enumerate}

\textit{\textcolor{grau}{Je 1P pro richtige Form}}

\vspace{15pt}

% ===============================================================
% AUFGABE 2 (mit Lösungen)
% ===============================================================

\begin{aufgabe}{2}
Setze die richtige Form des Verbs in Klammern ein. \punktebox{5}
\end{aufgabe}

\begin{enumerate}[label=\alph*)]
    \item Yo \lsg{vivo} en Rostock. (vivir)
    \item Tú \lsg{escribes} un mensaje. (escribir)
    \item Pedro \lsg{abre} la ventana. (abrir)
    \item Nosotros \lsg{vivimos} en Alemania. (vivir)
    \item Vosotros \lsg{escribís} cartas. (escribir)
\end{enumerate}

\textit{\textcolor{grau}{Je 1P pro richtige Form}}

\vspace{15pt}

% ===============================================================
% AUFGABE 3 (mit Lösungen)
% ===============================================================

\begin{aufgabe}{3}
Bilde Sätze mit den angegebenen Elementen. Konjugiere das Verb! \punktebox{4}
\end{aufgabe}

\begin{enumerate}[label=\alph*)]
    \item yo / leer / un libro / en casa\\
    \lsg{Yo leo un libro en casa.}

    \item mi hermano / correr / en el parque\\
    \lsg{Mi hermano corre en el parque.}

    \item nosotros / comer / pizza / con amigos\\
    \lsg{Nosotros comemos pizza con amigos.}

    \item tú / vivir / en España\\
    \lsg{Tú vives en España.}
\end{enumerate}

\textit{\textcolor{grau}{Je 1P pro korrekten Satz (Verb richtig konjugiert, Wortstellung akzeptabel)}}

\vspace{15pt}

% ===============================================================
% AUFGABE 4 (mit Lösungen)
% ===============================================================

\begin{aufgabe}{4}
Vervollständige die Tabelle und erkläre den Unterschied. \punktebox{6}
\end{aufgabe}

\begin{center}
\begin{tabular}{|l|c|c|c|}
\hline
\textbf{Person} & \textbf{hablar} (-ar) & \textbf{comer} (-er) & \textbf{vivir} (-ir) \\
\hline
yo & hablo & \lsg{como} & \lsg{vivo} \\
\hline
tú & hablas & \lsg{comes} & \lsg{vives} \\
\hline
él/ella & habla & \lsg{come} & \lsg{vive} \\
\hline
nosotros & hablamos & \lsg{comemos} & \lsg{vivimos} \\
\hline
vosotros & habláis & \lsg{coméis} & \lsg{vivís} \\
\hline
ellos & hablan & \lsg{comen} & \lsg{viven} \\
\hline
\end{tabular}
\end{center}

\textit{\textcolor{grau}{Tabelle: 4P (je 0,5P pro Zelle, max. 4P)}}

\vspace{1em}

\textbf{Erkläre:} Wo unterscheiden sich -er und -ir? (2P)

\lsg{Die Verben auf -er und -ir unterscheiden sich nur in der 1. und 2. Person Plural:}
\begin{itemize}
    \item \lsg{nosotros: -emos (comer) vs. -imos (vivir)}
    \item \lsg{vosotros: -éis (comer) vs. -ís (vivir)}
\end{itemize}

\textit{\textcolor{grau}{Je 1P für nosotros-Unterschied, 1P für vosotros-Unterschied}}

\vspace{20pt}
\hrule
\vspace{10pt}

\begin{flushright}
\textbf{Punkte:} \lsg{\maxpunkte} / \maxpunkte
\end{flushright}

\end{document}
